% =====================================================================
\chapter{序論}
\label{ch:intro}
% =====================================================================

\section{背景:量子多体系の性質を測るコスト}

量子計算機や量子シミュレータ(超伝導量子ビット、イオントラップ、冷却原子など)で多体の量子状態を用意できるようになり、その状態の性質を\emph{測定で取り出す}ことが中心的な課題になっている。量子力学の測定は確率的なので、ある物理量の期待値を精度 $\epsilon$ で知るには、おおよそ $1/\epsilon^2$ 回の測定を繰り返す必要がある。系が大きくなり、知りたい物理量が多くなるほど、この測定の回数(ショット数)と、測定のたびに回路を組み替える回数(測定の\emph{設定}の数)が実験の実質的なコストになる。

この問題に対する有力な手法が\emph{古典シャドウ}(classical shadow)\cite{huang2020}である。各量子ビットを独立にランダムな基底($X,Y,Z$ のいずれか)で測り、その結果から状態の「影」を古典計算機上に作る。1組の測定データから多数の局所的な物理量を同時に推定できることが利点である。一方で、観測量の台(作用する量子ビット)が $\lvert A\rvert$ 個のとき、推定の分散は $3^{\lvert A\rvert}$ に比例して増える。これは、ランダムに選んだ基底が観測量の台の上でたまたま全部一致したときにしか、その観測量の情報が得られないためである。

Vermersch らは、この古典シャドウに古典計算で作った近似状態を組み合わせる\emph{共通ランダム測定}(common randomized measurements, CRM)を提案した\cite{vermersch2024}。実験の影から「同じ基底で近似状態を測ったときの影」を差し引き、近似状態そのものを足し戻す。近似状態(以下 \emph{prior} と呼ぶ)が何であっても推定は偏らず、prior が良いほど分散が減る。統計学でいう\emph{共通乱数法}・\emph{制御変量法}の量子版である。

\section{本研究の問い}

CRM の効果は prior の良さで決まる。では、量子多体系、とくに強相関電子系の典型例であるハバード模型に対して、
\begin{enumerate}
  \item \textbf{どんな prior が、どの観測量で、どれだけ効くのか。}
  \item \textbf{それはなぜか} — prior のどの物理が、どの観測量の誤差に効いているのか。
\end{enumerate}
これが本研究の問いである。

元論文は、単一のパウリ文字列に対する CRM の分散を厳密に与えている(arXiv 版の付録 B.1 の式(25)、PRX Quantum 版の付録 C.1 の式(C7))。その式によれば、分散を決めるのは prior の\emph{大域的な忠実度}ではなく、その観測量に対する prior の外れ
\[
  \Delta=\Tr[P\rho]-\Tr[P\sigma]
\]
である(第\ref{ch:background}章)。したがって上の問いは、「ハバード模型の現実的な prior は、どの観測量でどれだけ外れるのか、それはなぜか」という\emph{物理の問い}に置き換わる。本研究の主な寄与はこの置き換えた問いに答えたことにある。

\section{本研究の寄与}

本研究で明らかにしたことを、元論文にあるものと区別して挙げる。

\begin{description}
  \item[元論文にあるもの(本研究は実装で確かめた)] 単一パウリ文字列の分散の厳密な式、その比として出る利得の式、損得の境目 $\lvert\Delta\rvert\le\lvert\Tr[P\rho]\rvert$、prior によらない不偏性。
  \item[本研究の寄与]
\end{description}
\begin{enumerate}
  \item \textbf{局所性}(第\ref{ch:locality}章):系を $L=8$ から $128$ サイト(256 量子ビット)に広げると prior の大域的な忠実度は崩れるが、局所観測量の利得はほとんど変わらないことを示した。prior の良さの指標としては、観測量の台に縮約したトレース距離が $\lvert\Delta\rvert$ の最良の上界で、半充填では等号が成り立つことも示した。
  \item \textbf{観測量ごとに prior の得意・不得意が逆転すること}(第\ref{ch:priors}章):平均場(非制限 Hartree--Fock, UHF)prior は電荷の量(二重占有に関わる量)では強結合で誤差が $(t/U)^2$ で消える一方、スピンの量では損をする。その原因を2つ(スピンの向きを選ぶこと、量子的な一重項の相関が足りないこと)に分け、それぞれの直し方と限界を示した。
  \item \textbf{prior の比較}(第\ref{ch:priors}章):スピン回転で平均した UHF(UHF-sym)、スピン射影 HF(PHF)、切断 MPS、変分 MPS、変分 MPS の回転平均を同じ厳密な参照状態で比べた。PHF の効果が $1/L$ で消えることを回転子の描像で定量的に説明した。変分 MPS を回転平均した prior が、厳密な状態を使わない prior の中で最も良いことを示した。
  \item \textbf{和の観測量の厳密な分散}(第\ref{ch:methods}章、付録\ref{app:sumvar}):二重占有のように複数のパウリ項の和で書かれる観測量について、CRM の分散を厳密に書き下した。利得は合計の $\Delta$ ではなく項ごとの外れで決まる。
  \item \textbf{2 次元とドープ系}(第\ref{ch:2d}章):最大 96 量子ビットの 2 次元シリンダーで、平均場 prior が MPS prior に代われる範囲と、ドープ系の電荷の模様で破綻する場所を調べた。
  \item \textbf{実用上の問題}(第\ref{ch:practical}章):ショットの配分、損をしない推定量(最適係数)、読み出し誤差に対する脆さと prior 側での補正、matchgate シャドウとの関係。
\end{enumerate}

\section{本論文の構成}

第\ref{ch:background}章でハバード模型・古典シャドウ・CRM を導入し、元論文の分散の式をまとめる。第\ref{ch:methods}章で本研究の手法(和の観測量の厳密な分散、prior の作り方、参照状態と数値実験)を述べる。第\ref{ch:locality}章から第\ref{ch:practical}章が結果である。第\ref{ch:conclusion}章で結論と今後の課題を述べる。付録に、半充填でサイトごとの電子数が厳密に 1 になる理由(粒子正孔対称性)、和の観測量の分散の導出、計算の再現方法をまとめた。
